\documentclass[12pt,a4paper]{article}
\usepackage[margin=2.5cm]{geometry}
\usepackage{amsmath,amssymb}
\usepackage{booktabs}
\usepackage{array}
\usepackage{enumitem}
\usepackage{parskip}
\usepackage{courier}
\usepackage[T1]{fontenc}
\usepackage{times}

\setlength{\parindent}{0pt}
\setlength{\parskip}{6pt}

% compact lists
\setlist[itemize]{noitemsep, topsep=2pt, leftmargin=1.4em, label={--}}

% monospace math blocks
\newcommand{\mline}[1]{\par\vspace{2pt}\hspace{1.5em}\texttt{#1}\par\vspace{2pt}}

\begin{document}

% ── Title ──────────────────────────────────────────────────────────────
\begin{center}
{\Large\textbf{IE 400 -- Term Project Report}}\\[4pt]
Spring 2026
\end{center}

\bigskip

% ═══════════════════════════════════════════════════════════════════════
\section*{\underline{Question 1: Museum Escape}}

\subsection*{Problem Definition}

Lara starts at \textbf{(1,1)} and must reach \textbf{(8,8)} on an $8\times8$ grid in the
minimum number of time steps. At each step she may move up, down, left, right, or stay in
place. Three security cameras patrol the museum; if Lara is ever in a camera-visible cell at
the same time step, she is caught.

\textbf{Obstacles:} (1,8), (2,3), (2,7), (3,6), (4,1), (4,2), (4,8), (5,6), (6,4),
(7,3), (7,4), (7,5)

\subsection*{Camera Movement Tables}

\textbf{Camera 1} -- period 4, horizontal patrol along row 3:

\begin{tabular}{cccc}
\toprule
$t$ & Position & Direction & Visible Cells \\
\midrule
0   & (3,3) & Right & (3,3),(3,4),(3,5) \\
1   & (3,4) & Right & (3,4),(3,5),(3,6) \\
2   & (3,5) & Left  & (3,5),(3,4),(3,3) \\
3   & (3,4) & Left  & (3,4),(3,3),(3,2) \\
4+  & (3,3) & Right & repeats from $t=0$ \\
\bottomrule
\end{tabular}

\medskip
\textbf{Camera 2} -- period 6, middle corridor:

\begin{tabular}{cccc}
\toprule
$t$ & Position & Direction & Visible Cells \\
\midrule
0  & (4,6) & Down  & (4,6),(5,6) \\
1  & (5,6) & Down  & (5,6),(6,6) \\
2  & (6,6) & Right & (6,6),(6,7) \\
3  & (6,7) & Left  & (6,7),(6,6) \\
4  & (6,6) & Up    & (6,6),(5,6) \\
5  & (5,6) & Up    & (5,6),(4,6) \\
6+ & (4,6) & Down  & repeats from $t=0$ \\
\bottomrule
\end{tabular}

\medskip
\textbf{Camera 3} -- period 4, exit area diagonal:

\begin{tabular}{cccc}
\toprule
$t$ & Position & Direction & Visible Cells \\
\midrule
0  & (7,7) & NE & (7,7),(6,8) \\
1  & (6,8) & SW & (6,8),(7,7),(8,6) \\
2  & (7,7) & SW & (7,7),(8,6) \\
3  & (8,6) & SW & (8,6),(7,7),(6,8) \\
4+ & (7,7) & NE & repeats from $t=0$ \\
\bottomrule
\end{tabular}

\medskip
The combined forbidden set at time $t$ is:
$\texttt{cam\_vis}(t) = C_1[t\bmod 4]\cup C_2[t\bmod 6]\cup C_3[t\bmod 4]$

% ── Task 1 ─────────────────────────────────────────────────────────────
\subsection*{Task 1: Integer Programming Formulation}

\textbf{Decision Variables}

$x[t,r,c]\in\{0,1\}$ -- equals 1 if Lara occupies cell $(r,c)$ at time step $t$, 0 otherwise.
Defined for all $t=0,1,\dots,T$ and all $(r,c)\in\text{VALID}$ (non-obstacle cells).

\textbf{Objective}

Minimize the total number of time steps to reach the exit:
\[
\text{Minimize } T
\]
$T^*$ is found first by time-expanded BFS (guaranteeing optimality); the IP is then solved at
that fixed horizon to confirm the solution and enforce task-specific constraints.

\textbf{Constraints}

\textbf{(C1) Initial position:}
\mline{x[0,1,1] = 1,\quad x[0,r,c] = 0 \; \forall\,(r,c)\neq(1,1)}

\textbf{(C2) One cell per time step:}
\[
\sum_{(r,c)\in\text{VALID}} x[t,r,c] = 1 \qquad \forall\,t
\]

\textbf{(C3) Valid movement} -- Lara may only move to an adjacent or same cell:
\[
x[t+1,r,c] \;\leq\; \sum_{(r',c')\in N(r,c)} x[t,r',c'] \qquad \forall\,t,\,(r,c)
\]
where $N(r,c)$ includes $(r,c)$ itself and the four orthogonal neighbours within the grid.

\textbf{(C4) Camera avoidance:}
\[
x[t,r,c] = 0 \qquad \forall\,t,\;\forall\,(r,c)\in\texttt{cam\_vis}(t),\;(r,c)\neq(8,8)
\]
The exit cell (8,8) is exempt -- Lara may safely reach it even if cameras see it.

\textbf{(C5) Exit reached:} $x[T,8,8]=1$

\textbf{(C6) Binary domain:} $x[t,r,c]\in\{0,1\}$

% ── Task 2 ─────────────────────────────────────────────────────────────
\subsection*{Task 2: Base Case}

A time-expanded BFS is run first over the state space (row, col, $t$) to find the minimum
horizon $T^*$. BFS is optimal here because each step costs exactly 1. The Gurobi IP is then
solved at $T=T^*$ to confirm the solution. A \texttt{validate()} function checks: no obstacle
entered, no camera-visible cell at the matching step, no jump larger than one cell, correct
start and end.

% ── Task 3 ─────────────────────────────────────────────────────────────
\subsection*{Task 3: Booby Trap}

Additional rule: if Lara visits \textbf{(5,1)}, she may never visit \textbf{(8,3)} afterwards.
This is an \emph{if-then} constraint modelled with a persistent binary variable:

$y[t]\in\{0,1\}$ -- equals 1 if Lara has visited (5,1) at or before time step $t$.

\mline{y[t]  >=  x[t,5,1]                   for all t}
\mline{y[t]  >=  y[t-1]                      for t >= 1  (monotone -- stays triggered)}
\mline{x[t,8,3]  <=  1 - y[t-1]             for t >= 1  (block (8,3) after trap)}

If no feasible solution exists at $T^*$, the code tries $T^*+1,\dots,T^*+9$.

% ── Task 4 ─────────────────────────────────────────────────────────────
\subsection*{Task 4: Exit Lock ($t < 18$)}

The cells $\{(6,7),(7,6),(7,7),(8,6),(6,8)\}$ are forbidden before $t=18$:
\[
x[t,r,c]=0 \qquad \forall\,t<18,\;\forall\,(r,c)\in\text{LOCKED}
\]
Both the BFS (via \texttt{bfs\_block4}) and the IP (via \texttt{extra4}) enforce this.

% ── Task 5 ─────────────────────────────────────────────────────────────
\subsection*{Task 5: Ghost Checkpoint ($t=3$)}

Lara must be at \textbf{(4,4)} exactly at $t=3$:
\[
x[3,4,4]=1
\]
BFS is guided with \texttt{forced\_at = \{3:(4,4)\}} and the IP adds this equality constraint.
If infeasible at $T^*$, horizons up to $T^*+14$ are tried.

% ── Task 6 ─────────────────────────────────────────────────────────────
\subsection*{Task 6: Branch and Bound (Base Case)}

Branch and Bound (B\&B) solves the IP by partitioning the feasible region into subproblems
(\emph{branching}) and using LP relaxation bounds to discard subproblems that cannot improve
the best known solution (\emph{bounding}).

\textbf{LP Relaxation.}
Relax $x[t,r,c]\in\{0,1\}$ to $x[t,r,c]\in[0,1]$. The LP optimum provides a lower bound.
Since $T$ is fixed at $T^*$ by BFS, the IP is a feasibility problem -- B\&B finds the first
$\{0,1\}$ assignment satisfying all constraints.

\textbf{Branching strategy.}
At each node, solve the LP relaxation. If any variable $x[t,r,c]$ is fractional (e.g.\ 0.5),
branch by creating two children:
\begin{itemize}
  \item Branch 0: add $x[t,r,c]=0$
  \item Branch 1: add $x[t,r,c]=1$
\end{itemize}
Variable selected: most fractional (closest to 0.5). Node selected: best-first.

\textbf{Pruning.}
A node is pruned when (i) its LP relaxation is infeasible, or (ii) it cannot improve the best
known bound.

The museum IP is highly constrained -- uniqueness, movement, and camera constraints push the LP
relaxation very close to integrality. In practice B\&B terminates within 1--2 branching levels
and recovers the same path found by BFS.

\begin{tabular}{lll}
\toprule
Node & Action & Result \\
\midrule
Root     & Solve LP relaxation ($x\in[0,1]$) & Near-integer or integer solution \\
Branch 0 & Fix most-fractional $=0$           & Prune if LP infeasible \\
Branch 1 & Fix most-fractional $=1$           & Update incumbent if integer \\
Leaf     & All vars $\in\{0,1\}$              & Optimal path confirmed \\
\bottomrule
\end{tabular}

\newpage

% ═══════════════════════════════════════════════════════════════════════
\section*{\underline{Question 2: CCI Production Planning}}

\subsection*{Problem Definition}

Coca-Cola \.{I}\c{s}ecek (CCI) plans production at two plants -- \textbf{Ankara (A)} and
\textbf{Istanbul (I)} -- over 6 months. Minimize total cost (inventory holding, transportation,
wages, hiring, layoffs, bottling mode switching) subject to demand, capacity, safety stock,
workforce, and cross-shipping constraints.

\subsection*{Data}

\textbf{Monthly demand (million liters):}

\begin{tabular}{lcccccc}
\toprule
Region & M1 & M2 & M3 & M4 & M5 & M6 \\
\midrule
Western (W)          & 9 & 9 & 14 & 12 & 12 & 13 \\
Central \& Eastern (CE) & 7 & 7 &  8 &  9 &  8 &  9 \\
\bottomrule
\end{tabular}

\medskip
\textbf{Plant parameters:}

\begin{tabular}{lll}
\toprule
Parameter & Ankara (A) & Istanbul (I) \\
\midrule
Bottling rate -- 2L & 6,000 bot/hr = 12,000 L/hr & 9,000 bot/hr = 18,000 L/hr \\
Bottling rate -- 1L & 10,000 bot/hr = 10,000 L/hr & 15,000 bot/hr = 15,000 L/hr \\
Available hours/month & 720 hr & 720 hr \\
Initial workforce & 60 emp & 100 emp \\
Primary region & Central \& Eastern & Western \\
\bottomrule
\end{tabular}

\medskip
\textbf{Costs:} wage 1,200~TL/emp$\cdot$month; hire 1,500~TL/emp; layoff 2,000~TL/emp;
inventory 0.05~TL/L$\cdot$month (= 50,000~TL/ML); cross-region transport 0.10~TL/L
(= 100,000~TL/ML); bottling mode switch 20,000~TL/plant$\cdot$month.

\subsection*{Decision Variables}

All production/inventory quantities in million liters (ML). $m\in\{1,\dots,6\}$.
\begin{itemize}
  \item $pA1[m],\,pA2[m]\geq0$ -- 1L / 2L production at Ankara (ML)
  \item $pI1[m],\,pI2[m]\geq0$ -- 1L / 2L production at Istanbul (ML)
  \item $sAW[m],\,sACE[m]\geq0$ -- Ankara supply to Western / C\&E (ML)
  \item $sIW[m],\,sICE[m]\geq0$ -- Istanbul supply to Western / C\&E (ML)
  \item $invA[m],\,invI[m]\geq0$ -- end-of-month inventory at Ankara / Istanbul (ML)
  \item $wA[m],\,wI[m]\in\mathbb{Z}_{\geq0}$ -- workforce; $hA,hI\geq0$ hired; $lA,lI\geq0$ laid off
  \item $bA1,bA2,bI1,bI2[m]\in\{0,1\}$ -- bottling mode indicators (1L/2L active per plant-month)
  \item $zA[m],\,zI[m]\in\{0,1\}$ -- switching indicators (1 if both modes used at a plant in month $m$)
\end{itemize}

\subsection*{Objective Function}

\[
\min\; Z = C_{\text{inv}}+C_{\text{trans}}+C_{\text{wage}}+C_{\text{hire}}+C_{\text{lay}}+C_{\text{sw}}
\]
\begin{align*}
C_{\text{inv}}   &= 50{,}000  \sum_m (invA[m]+invI[m]) \\
C_{\text{trans}} &= 100{,}000 \sum_m (sAW[m]+sICE[m]) \\
C_{\text{wage}}  &= 1{,}200   \sum_m (wA[m]+wI[m]) \\
C_{\text{hire}}  &= 1{,}500   \sum_m (hA[m]+hI[m]) \\
C_{\text{lay}}   &= 2{,}000   \sum_m (lA[m]+lI[m]) \\
C_{\text{sw}}    &= 20{,}000  \sum_m (zA[m]+zI[m])
\end{align*}

\subsection*{Constraints}

\textbf{(D1) Bottling capacity} -- total hours used $\leq$ 720:
\[
\frac{pA2[m]\cdot10^6}{12{,}000}+\frac{pA1[m]\cdot10^6}{10{,}000}\leq720 \quad
\frac{pI2[m]\cdot10^6}{18{,}000}+\frac{pI1[m]\cdot10^6}{15{,}000}\leq720 \qquad \forall\,m
\]

\textbf{(D2) Packaging capacity} -- total ML $\leq$ workforce $\times$ 0.1:
\[
pA1[m]+pA2[m]\leq0.1\,wA[m], \qquad pI1[m]+pI2[m]\leq0.1\,wI[m] \qquad \forall\,m
\]

\textbf{(D3) Workforce balance} ($wA[0]=60$, $wI[0]=100$):
\[
wA[m]=wA[m-1]+hA[m]-lA[m], \qquad wI[m]=wI[m-1]+hI[m]-lI[m] \qquad \forall\,m
\]

\textbf{(D4) Workforce smoothing} ($\Delta=10$):
\[
|wA[m]-wA[m-1]|\leq10, \qquad |wI[m]-wI[m-1]|\leq10 \qquad \forall\,m
\]
(Linearised as two constraints each: $wA[m]-wA[m-1]\leq10$ and $wA[m-1]-wA[m]\leq10$.)

\textbf{(D5) Inventory balance} (initial inventories $=0$):
\begin{align*}
invA[m]&=invA[m-1]+pA1[m]+pA2[m]-sACE[m]-sAW[m]\\
invI[m]&=invI[m-1]+pI1[m]+pI2[m]-sIW[m]-sICE[m]
\end{align*}

\textbf{(D6) Demand fulfillment:}
\[
sACE[m]+sICE[m]\geq D_{\text{CE}}[m], \qquad sIW[m]+sAW[m]\geq D_W[m] \qquad \forall\,m
\]

\textbf{(D7) Cross-shipping cap} ($\alpha=25\%$):
\[
sAW[m]\leq0.25\,D_W[m], \qquad sICE[m]\leq0.25\,D_{\text{CE}}[m] \qquad \forall\,m
\]

\textbf{(D8) Safety stock} ($\beta=10\%$, months 1--5):
\[
invI[m]\geq0.10\,D_W[m+1], \qquad invA[m]\geq0.10\,D_{\text{CE}}[m+1] \qquad m=1,\dots,5
\]

\textbf{(D9) Zero ending inventory:}
\[
invA[6]=0, \qquad invI[6]=0
\]

\textbf{(D10) Bottling mode switching cost} (big-$M=50$ ML):
\[
pA1[m]\leq50\,bA1[m],\quad pA2[m]\leq50\,bA2[m],\quad
pI1[m]\leq50\,bI1[m],\quad pI2[m]\leq50\,bI2[m] \qquad \forall\,m
\]
\[
zA[m]\geq bA1[m]+bA2[m]-1, \qquad zI[m]\geq bI1[m]+bI2[m]-1 \qquad \forall\,m
\]
$zA[m]=1$ iff both bottle sizes are produced at Ankara in month $m$ (same logic for $zI$),
triggering the 20,000~TL switching cost.

\end{document}
